\documentclass[11pt]{article}

\usepackage[margin=1in]{geometry}
\usepackage{booktabs}
\usepackage{array}
\usepackage{graphicx}
\usepackage{hyperref}

\newcommand{\code}[1]{\texttt{#1}}

\title{C5G7 2-D CSG Benchmark Report}
\author{Detran example suite}
\date{\today}

\begin{document}
\maketitle

\section*{Purpose}

This report is a self-contained draft for comparing Detran C5G7 calculations
against the published two-dimensional C5G7 benchmark reference values.  The
reference eigenvalue used by the example suite is
\[
  k_{\mathrm{eff}}^{\mathrm{ref}} = 1.186572165 ,
\]
stored in \code{c5g7/reference.py}.  The original C5G7 reference solution is a
Monte Carlo benchmark reference; in this report we use ``MCNP reference'' as
shorthand for that published reference data.

\section*{CSG Geometry}

The CSG variant is generated by \code{c5g7/csg\_c5g7.py}.  It preserves the
existing C5G7 material IDs and pin maps:

\begin{center}
\begin{tabular}{cl}
\toprule
Material ID & Description \\
\midrule
0 & UO2 fuel \\
1 & MOX 4.3\% fuel \\
2 & MOX 7.0\% fuel \\
3 & MOX 8.7\% fuel \\
4 & Guide tube \\
5 & Fission chamber \\
6 & Moderator \\
\bottomrule
\end{tabular}
\end{center}

The main pin-level CSG discretization controls are:

\begin{itemize}
  \item \code{fuel\_radial\_divisions}: radial divisions in the fuel or central circular material.
  \item \code{moderator\_radial\_divisions}: radial divisions in the moderator outside the central radius and clipped by the pin square.
  \item \code{azimuthal\_slices}: angular sectors.  The default is 8, which gives sector boundaries every 45 degrees, including the pin diagonals.
\end{itemize}

The default CSG pin therefore has one central radial zone, one moderator radial
zone, and eight azimuthal sectors.  Fuel, MOX, guide-tube, and fission-chamber
pins have 16 CSG regions under the default; pure moderator pins have 8 CSG
regions.  These CSG regions are also the flat-source regions used by the
current native CSG-MOC bridge.

\begin{figure}[htbp]
\centering
\IfFileExists{results/figures/csg_geometry_material_maps.pdf}{%
  \includegraphics[width=\textwidth]{results/figures/csg_geometry_material_maps.pdf}%
}{}
\caption{Detran-sampled CSG material maps for the UO2 pin, UO2 assembly, and
quarter-core geometry.  The figure is generated from \code{CSGModel2D.sample\_slice},
so region classification comes from the Detran CSG backend rather than from a
separate plotting reconstruction.}
\end{figure}

\begin{figure}[htbp]
\centering
\IfFileExists{results/figures/csg_flat_source_regions.pdf}{%
  \includegraphics[width=\textwidth]{results/figures/csg_flat_source_regions.pdf}%
}{}
\caption{Representative Detran-sampled flat-source region maps.  Colors are
hashed from \code{region\_id} values returned by \code{CSGModel2D.sample\_slice};
black edges mark sampled region transitions, and white edges mark pin-cell
boundaries in the assembly zooms.}
\end{figure}

\section*{Run Cases}

Each driver writes a JSON result and a LaTeX row fragment under
\code{csg\_results/}.  The report table below loads those fragments when they
exist.

\begin{verbatim}
python3 -m c5g7.run_case_csg_pin --track --solve-csg-moc
python3 -m c5g7.run_case_csg_assembly --track --solve-csg-moc
python3 -m c5g7.run_case_csg_core --track --solve-csg-moc
\end{verbatim}

This path builds a product-quadrature CSG \code{TrackCollection}, constructs a
region-indexed synthetic mesh with one cell per CSG region, and forwards the
explicit track product into the Detran MOC eigenvalue manager.  It therefore
uses the existing MOC machinery for boundary conditions, fission source updates,
Krylov/vector construction, and sweeps, while replacing mesh-derived tracks with
CSG-derived region tracks.

The current default CSG-MOC artifact set uses one polar angle per octant, four
azimuthal angles per octant, and a maximum track spacing of 0.25 cm.  It now
contains direct Detran eigenvalue rows for the UO2 pin, UO2 assembly, and
quarter-core cases.  These rows are first-pass transport checks and refinement
baselines, not final reference-quality calculations.

For comparison, the CSG drivers can still include the older staircase mesh-MOC
solve:

\begin{verbatim}
python3 -m c5g7.run_case_csg_pin --track --solve-mesh-moc
\end{verbatim}

The CSG bridge currently estimates per-region areas from the CSG track product
and scales them to the analytic domain area before creating the synthetic mesh.
This is sufficient for smoke and refinement studies, but exact analytic region
areas remain follow-on work.  Quarter-core pin powers are reduced by summing
\(\Sigma_f \phi V\) over the CSG regions belonging to each fuel pin and
normalizing the active-pin sum to the MCNP reference pin-power sum.

\section*{SN/DD Discretization}

The first concrete SN variant uses the mesh-backed two-dimensional Detran
transport path with diamond difference, PETSc/SLEPc-enabled eigenvalue solvers,
and MG-CMDSA or CMFD acceleration options.  The artifact drivers are:

\begin{verbatim}
python3 -m c5g7.run_case_sn_pin
python3 -m c5g7.run_case_sn_assembly
python3 -m c5g7.run_case_sn_core
\end{verbatim}

The default solver setting is:
\begin{verbatim}
--eigen-backend slepc-gd --acceleration mgcmdsa --linear-backend petsc
\end{verbatim}
The same drivers can also run
\code{gd}, \code{slepc-krylovschur}, \code{cmfd}, or \code{power} for fallback
and comparison studies.

The current comparison rows use the same coarse angular count as the first-pass
CSG-MOC rows: one polar angle per octant and four azimuthal angles per octant
with \code{u-u} product quadrature.  The SLEPc-GD/MGCMDSA assembly path was
substantially more expensive in this smoke configuration, so these rows use
power iteration with one gauss-seidel outer sweep per eigen step, matching the
legacy C5G7 SN example style.

\subsection*{Automated 7x7 Pin Mesh}

The automated 7x7 pin mesh preserves the historical D4-symmetric 13-cell
central-material mask,
\[
\begin{array}{ccccccc}
0&0&0&0&0&0&0\\
0&0&0&1&0&0&0\\
0&0&1&1&1&0&0\\
0&1&1&1&1&1&0\\
0&0&1&1&1&0&0\\
0&0&0&1&0&0&0\\
0&0&0&0&0&0&0
\end{array}
\]
but computes the mesh widths automatically.  For pin pitch \(p\), radius \(r\),
central square width \(L\), and one-step arm thickness \(\delta\), the symmetric
width vector is
\[
  [m,\delta,L/4,L/2,L/4,\delta,m],
  \qquad
  m = p/2 - L/2 - \delta .
\]
The central-material area of this mask is the square plus four arms,
\[
  A_h = L^2 + 4\left(\frac{L}{2}\delta\right)
      = L^2 + 2L\delta .
\]
Volume preservation imposes \(A_h = \pi r^2\).  With \(k=L/r\),
\[
  \frac{\delta}{r} = \frac{\pi-k^2}{2k},
  \qquad
  m = \frac{p}{2} - \frac{kr}{2}
      - \frac{(\pi-k^2)r}{2k}.
\]
The remaining scalar \(k\) is selected by minimizing the normalized symmetric
difference between the staircase central region \(S(k)\) and the true circle
\(C\),
\[
  J(k) = \frac{|S(k)\triangle C|}{\pi r^2},
\]
using deterministic Gauss-Legendre rectangle-circle overlap integrals.  This
gives \(k \approx 1.61008\), very close to the legacy constant \(1.60496875\).
The resulting C5G7 widths are approximately
\[
 [0.10317444,\; 0.09210462,\; 0.21736047,\; 0.43472094,\;
  0.21736047,\; 0.09210462,\; 0.10317444]\ \mathrm{cm}.
\]

The implemented sanity checks enforce positive symmetric widths, exact pitch
closure, 13 central cells, central area preservation to roundoff, moderator area
preservation to roundoff, and a symmetric-difference fraction no worse than the
legacy 7x7 constants.  For the C5G7 pin dimensions, the optimized symmetric
difference fraction is about \(7.624\times10^{-2}\), compared with
\(7.632\times10^{-2}\) for the legacy constants.

\section*{Results}

\begin{center}
\scriptsize
\setlength{\tabcolsep}{2.5pt}
\begin{tabular}{lrrrrrrlrrr}
\toprule
Case & Fuel $r$ & Mod. $r$ & Az. & Regions & Tracks & Segments
     & Backend & Ref. $k$ & Detran $k$ & pcm \\
\midrule
\IfFileExists{csg_results/csg_pin_0_row.tex}{\input{csg_results/csg_pin_0_row.tex}\\}{}
\IfFileExists{csg_results/csg_assembly_0_row.tex}{\input{csg_results/csg_assembly_0_row.tex}\\}{}
\IfFileExists{csg_results/csg_core_0_row.tex}{\input{csg_results/csg_core_0_row.tex}\\}{}
\bottomrule
\end{tabular}
\end{center}

\subsection*{Few-pcm UO2 Pin-Cell Accuracy Runs}

The UO2 pin-cell reference in \code{c5g7/reference.py} is the historical
Detran 7x7 volume-preserving pin-cell regression value.  The following rows
therefore use 7x7 volume-preserving pin meshes and tighter angular/solver
settings rather than the earlier smoke settings.  Both rows are direct solves;
the pcm column is \((k-k_{\mathrm{ref}})10^5\) with
\(k_{\mathrm{ref}}=1.325939322\).

\input{results/uo2_pin_accuracy_table.tex}

\subsection*{Centered NxN UO2 Pin-Cell Diagnostic Sweeps}

The centered Cartesian UO2 pin-cell sweeps provide the first common spatial
refinement ladder for mesh-backed native MOC and SN, but they are diagnostic
relative to the 7x7 regression reference above.  Each row uses reflected pin
boundaries, UO2 pin index 0, a centered \(N\times N\) material assignment, and
density-corrected macroscopic cross sections so the true fuel and moderator
areas are preserved even when the cell-center material mask cuts the circle.
The common angular setting is \code{u-gl} product quadrature with two polar
angles and four azimuthal angles per octant.  For the MOC rows, \(h\) is the
track-spacing value selected from the minimum cell width; for the SN rows, \(h\)
is the minimum mesh-cell width.

\input{results/mesh_sweep_uo2_pin_table.tex}

These centered-mesh rows are raw direct solves, not extrapolated reference
estimates.  The MOC and SN-DD sequences agree near
\(k_{\mathrm{eff}}\approx 1.3238\) on the finest completed shared mesh, roughly
\(-220\) pcm against the 7x7 UO2 pin-cell reference.  The gap is therefore not
a few-pcm solver-convergence issue; it reflects that the centered
density-corrected mesh is a different approximation from the volume-preserving
7x7 reference mesh.  The SN-SC sequence is included because it was run on the
same ladder, but its current flattened power-iteration settings do not show an
asymptotic spatial trend and should be treated as diagnostic.

\subsection*{Core Pin-Power Errors}

\begin{center}
\scriptsize
\setlength{\tabcolsep}{4pt}
\begin{tabular}{lrrrrlrr}
\toprule
Method & Active pins & RMS [\%] & Mean abs. [\%] & Max abs. [\%]
       & Max pin $(i,j)$ & Ref. & Detran \\
\midrule
\IfFileExists{csg_results/csg_core_0_pin_error_row.tex}{\input{csg_results/csg_core_0_pin_error_row.tex}\\}{}
\bottomrule
\end{tabular}
\end{center}

\begin{figure}[htbp]
\centering
\IfFileExists{results/figures/csg_core_pin_power_errors.pdf}{%
  \includegraphics[width=\textwidth]{results/figures/csg_core_pin_power_errors.pdf}%
}{}
\caption{Quarter-core pin-power reference map, CSG-MOC computed map, and
relative error map.  Reflector pins are masked.  Assembly boundaries are drawn
every 17 pins.}
\end{figure}

\begin{center}
\tiny
\setlength{\tabcolsep}{1.5pt}
\begin{tabular}{llrllllllrrrr}
\toprule
Case & Mesh & Cells & Eq. & Quad. & Pol. & Az. & Backend & Accel.
     & Ref. $k$ & Detran $k$ & pcm & Time [s] \\
\midrule
\IfFileExists{results/sn/sn_pin_0_optimized_row.tex}{\input{results/sn/sn_pin_0_optimized_row.tex}\\}{}
\IfFileExists{results/sn/sn_assembly_0_optimized_row.tex}{\input{results/sn/sn_assembly_0_optimized_row.tex}\\}{}
\IfFileExists{results/sn/sn_core_0_optimized_row.tex}{\input{results/sn/sn_core_0_optimized_row.tex}\\}{}
\bottomrule
\end{tabular}
\end{center}

\section*{Interpretation}

The CSG geometry is now the exact circular-pin representation needed for a
deterministic C5G7 study, with explicit knobs for pin radial and azimuthal
subdivision.  Native CSG-MOC transport now reaches the Detran eigenvalue manager
through a region-track bridge, so CSG rows can contain direct Detran
$k_{\mathrm{eff}}$ values rather than only mesh-MOC comparisons.  The current
pin-power error map shows that this is not yet a benchmark-quality setting:
with the default one-radial-zone, eight-sector pin CSG and 0.25 cm MOC spacing,
the quarter-core RMS pin-power error is still in the double digits.  The
remaining CSG work is accuracy-oriented: exact analytic region volumes,
refinement studies over CSG subdivision, track spacing, and angular quadrature,
and then sequence/extrapolation studies once the raw ladder is credible.

The SN/DD path is now concrete and reproducible.  It can run the legacy
7x7 mesh or the automated 7x7 recipe from the same artifact-writing interface,
and the solver metadata distinguishes SLEPc GD, native GD, CMFD, and power
iteration comparisons.

\end{document}
